\chapter{The Fractal Architecture}
\label{app:fractal_architecture}
\addtocontents{toc}{\protect\setcounter{tocdepth}{0}}

\textbf{Subject:} Vertical Symmetry of Action ($h$) and Gravitational Coupling ($G$) Projected from Substrate to Soft Floor to Matter.

\section{The Fractal Hypothesis (Source \& Projection)}
In \hyperref[app:emergence_h_G]{Appendix \ref*{app:emergence_h_G}}, we derived the fundamental constants $G$ and $h$ using the specific mechanical parameters of the Soft Floor ($J_D$). In this appendix, we demonstrate that these constants are not unique to the vacuum; they are geometric projections from the Substrate that re-emerge at specific observable scales.

\textbf{The Generalized Hypothesis:} The observable constants are scale-invariant. They emerge fractally from the Universal Torsional Invariant ($J_i = \rho_i r_i^4$)---the strict mechanical property that dictates both internal geometric action and external substrate coupling across all phases of the universe.

Tension Field Theory (TFT) posits that the universe operates on a scale-invariant architecture. We define the constants based on their layer of origin:
\begin{itemize}
    \item \textbf{Layer S (Substrate):} The Hard Floor SC lattice ($\ell_S$). The absolute source of latent $\hbar$ and $G$.
    \item \textbf{Layer D (Soft Floor):} The vacuum FCC lattice ($r_D$). Where $h$ and $G$ project to form Spacetime.
    \item \textbf{Layer N (Nucleon):} The matter vortex ($r_N$). Where $h$ and $G$ project to form Mass.
\end{itemize}

\textbf{The Fundamental Fractal Key:} The projection is strictly governed by the conservation of the Universal Torsional Invariant ($J_i = \rho_i r_i^4$) across scales.
\begin{itemize}
    \item \textbf{Action:} $h_{layer} \propto (\text{Geometric Factor}) \cdot (J_i \cdot c_E)$
    \item \textbf{Gravity:} $G_{layer} \propto \dfrac{1}{J_i}$
\end{itemize}

\section{Thread A: The Projection of Action ($h$)}
The action of existence is defined by the Torsional Invariant of the node ($J_i$) modulated by its specific phase geometry.

\subsection{Level I: The Source (Layer S)}
\begin{itemize}
    \item \textbf{State:} Static substrate (pure linear radius).
    \item \textbf{Geometry:} Radial vector ($\text{factor} = 1.0$).
\end{itemize}
\textbf{Calculation:}
\begin{equation}
    h_S \approx \rho_S c_E \ell_S^4 \approx 1.0544 \times 10^{-34} \text{ J}\cdot\text{s} \to \hbar
\end{equation}
\textit{(The absolute Substrate dictates the reduced Planck constant, $\hbar$).}

\subsection{Level II: The Projection to Vacuum (Layer D)}
\begin{itemize}
    \item \textbf{State:} Phase-locked solid (sphere).
    \item \textbf{Scale:} $r_D \approx 7.8900 \times 10^{-27} \text{ m}$.
    \item \textbf{Density:} $\rho_D \approx 1.3615 \times 10^{62} \text{ kg/m}^3$.
    \item \textbf{Geometry:} Solid sphere ($\text{factor} = \frac{4}{3}\pi$).
\end{itemize}
\textbf{Calculation:}
\begin{equation}
    h_D = \frac{4}{3}\pi \rho_D c_E r_D^4 \approx 6.6257 \times 10^{-34} \text{ J}\cdot\text{s} \to h
\end{equation}

\subsection{Level III: The Projection to Matter (Layer N)}
\begin{itemize}
    \item \textbf{State:} Topological resonance (phase-locked knot).
    \item \textbf{Scale:} $r_N \approx 0.88000 \times 10^{-15} \text{ m}$ (neutron mechanical radius).
    \item \textbf{Density:} $\rho_N \approx 5.8658 \times 10^{17} \text{ kg/m}^3$.
    \item \textbf{Geometry:} Cylindrical vortex ($\text{factor} = 2\pi$).
\end{itemize}
\textbf{Calculation:}
\begin{equation}
    h_N = 2\pi \rho_N c_E r_N^4 \approx 6.6261 \times 10^{-34} \text{ J}\cdot\text{s} \to h
\end{equation}

\textbf{Conclusion:} Matter is a macroscopic resonance of the vacuum. The geometric factor increases ($\frac{4}{3}\pi \to 2\pi$) exactly as the Torsional Invariant decreases, perfectly conserving the macroscopic action $h$.

\section{Thread B: The Projection of the Gravitational Constant ($G$)}
The gravitational constant $G$ is the coupling efficiency between the layer's Torsional Invariant ($J_i$) and the absolute vacuum Substrate Area ($A_S = \ell_S^2 \approx 2.6123 \times 10^{-70} \text{ m}^2$).

\textbf{Master Equation:}
\begin{equation}
    G_i = (\text{Form Factor}) \cdot \left(\frac{c_E^2 A_S}{J_i}\right)
\end{equation}

\subsection{Level I: The Source (Layer S)}
At the Substrate scale, the node couples to itself. The form factor is $1.0$.
\begin{itemize}
    \item \textbf{Torsional Invariant:} $J_S = \rho_S \ell_S^4 \approx 3.5172 \times 10^{-43} \text{ kg}\cdot\text{m}$.
\end{itemize}
\textbf{Calculation:}
\begin{equation}
    G_S = 1.0 \cdot \left(\frac{c_E^2 \ell_S^2}{\rho_S \ell_S^4}\right) \approx 6.6748 \times 10^{-11} \text{ m}^3/(\text{kg}\cdot\text{s}^2) \equiv G
\end{equation}

\subsection{Level II: The Projection to Vacuum (Layer D)}
\begin{itemize}
    \item \textbf{State:} Spherical coupling.
    \item \textbf{Torsional Invariant:} $J_D = \rho_D r_D^4 \approx 5.2763 \times 10^{-43} \text{ kg}\cdot\text{m}$.
    \item \textbf{Form Factor:} $1.5$ (spherical surface coupling).
\end{itemize}
\textbf{Calculation:}
\begin{equation}
    G_D = 1.5 \cdot \left(\frac{c_E^2 A_S}{J_D}\right) \approx 6.6742 \times 10^{-11} \text{ m}^3/(\text{kg}\cdot\text{s}^2) \equiv G
\end{equation}

\subsection{Level III: The Projection to Matter (Layer N)}
\begin{itemize}
    \item \textbf{State:} Cylindrical coupling (vortex).
    \item \textbf{Torsional Invariant:} $J_N = \rho_N r_N^4 \approx 3.5177 \times 10^{-43} \text{ kg}\cdot\text{m}$. \textit{(Note: $J_N \approx J_D / 1.5$. The nucleon is exactly $1.5\times$ softer than the local spacetime lattice).}
    \item \textbf{Form Factor:} $1.0$ (linear/cylindrical coupling).
\end{itemize}
\textbf{Calculation:}
\begin{equation}
    G_N = 1.0 \cdot \left(\frac{c_E^2 A_S}{J_N}\right) \approx 6.6738 \times 10^{-11} \text{ m}^3/(\text{kg}\cdot\text{s}^2) \equiv G
\end{equation}

\textbf{The Impedance Match ($J_N \approx J_S$):}
\begin{itemize}
    \item $J_S \approx 3.5172 \times 10^{-43} \text{ kg}\cdot\text{m}$
    \item $J_N \approx 3.5177 \times 10^{-43} \text{ kg}\cdot\text{m}$
\end{itemize}

\textbf{Conclusion:} The Torsional Invariant of a nucleon is virtually identical to the Torsional Invariant of the absolute Substrate. This mathematically reveals that matter is a superconductive channel that tunnels through the Soft Floor vacuum tension to connect directly to the Hard Floor.

\section{Conclusion: The Unified Architecture}
The universe is a vertical fractal defined by geometric phase-locking:
\begin{itemize}
    \item \textbf{The Hard Floor ($10^{-35} \text{ m}$):} Sets the absolute grain ($\hbar$, $G$).
    \item \textbf{The Soft Floor ($10^{-26} \text{ m}$):} Splits reality into a solid node (source of $h$ via $\frac{4}{3}\pi$ and modifier of $G$ via spherical form factor $1.5$) and interstitial condensate (the source of kinematic inertia).
    \item \textbf{The Nucleon ($10^{-15} \text{ m}$):} Is the macroscopic realization of this split. It conserves $h$ by shifting geometry to a vortex ($2\pi$) and conserves $G$ by lowering its stiffness to perfectly match the Hard Floor limit ($J_N = J_S$).
\end{itemize}

There are not three physics. There is one fundamental torsional geometry ($J_i = \rho_i r_i^4$) echoing through the inflating topology of spacetime.

\vspace{0.5cm} % Forces a clean, half-centimeter vertical drop
\noindent \textbf{Note on the Nucleon Radius Anomaly:} Current experimental physics reports a discrepancy in the proton/neutron radius between $0.84 \text{ fm}$ (charge radius) and $0.88 \text{ fm}$ (CODATA/scattering). The TFT fractal architecture resolves this anomaly:
\begin{itemize}
    \item $0.88 \text{ fm}$ ($r_N$): Is the mechanical radius (the vortex loop) required to conserve $G$ and $h$ by strictly matching the Substrate impedance ($J_N = J_S$).
    \item $0.84 \text{ fm}$: Is the charge distribution radius (the electric field halo), which naturally sits slightly inside the mechanical boundary.
\end{itemize}
TFT predicts that future experiments measuring gravitational or inertial properties will converge absolutely on $0.88 \text{ fm}$.